\subsection{Tensor product of point distributions}

13.4.1. Let \(\mathbf{X}_1\) and \(\mathbf{X}_2\) be two manifolds of class \(\mathbf{C}^r\) , let \(x_1 \in \mathbf{X}_1\) , \(x_2 \in \mathbf{X}_2\) and \(t_1 \in \mathrm{T}_{x_1}^{(k_1)}(\mathbf{X}_1)\) , \(t_2 \in \mathrm{T}_{x_2}^{(k_2)}(\mathbf{X}_2)\) with \(k_1 + k_2 \leqslant r\) . Set \(\mathbf{X} = \mathbf{X}_1 \times \mathbf{X}_2\) , \(x = (x_1, x_2)\) and \(k = k_{1} + k_{2}\) . For \(i = 1, 2\) let \(c_{i} = (U_{i}, \varphi_{i}, E_{i})\) be a chart of \(X_{i}\) centered at \(x_{i}\) and denote by \(\tilde{t}_{i}\) the element of \(TS(E_{i})\) such that \(\theta_{c_i}(\tilde{t}_{i}) = t_{i}\) , cf. 13.2.1. Let \(\sigma\) be the canonical isomorphism of \(TS(E_{1}) \otimes TS(E_{2})\) onto \(TS(E_{1} \times E_{2})\) , cf. A, IV, § 5, no. 5. The element \(\sigma(\tilde{t}_{1} \otimes \tilde{t}_{2})\) is the symmetric product of \(\tilde{t}_{1}\) and \(\tilde{t}_{2}\) (identified with elements of \(TS(E_{1} \times E_{2})\) by means of the canonical injections \(TS(E_{i}) \to TS(E_{1} \times E_{2})\) ); it belongs to \(TS^{(k)}(E_{1} \times E_{2})\) . Set \(c = c_{1} \times c_{2}\) ; it is a chart of X centered at \(x\) . The image under \(\theta_{c}\) of \(\sigma(\tilde{t}_{1} \otimes \tilde{t}_{2})\) is an element of \(T_{x}^{(k)}(X)\) , which does not depend on the choice of charts \(c_{i}\) . It is called the direct product, or the symmetric tensor product (or even simply the tensor product) of \(t_{1}\) and \(t_{2}\) and is denoted \(t_{1} \times t_{2}\) or \(t_{1} \otimes t_{2}\) .

One has \(\varepsilon_{x_1} \otimes \varepsilon_{x_2} = \varepsilon_x\) . The constant term of \(t_1 \otimes t_2\) is the product of the constant terms of \(t_1\) and of \(t_2\) .

The isomorphism \((y_{1}, y_{2}) \mapsto (y_{2}, y_{1})\) of \(\mathbf{X}_{1} \times \mathbf{X}_{2}\) onto \(\mathbf{X}_{2} \times \mathbf{X}_{1}\) transforms \(t_{1} \otimes t_{2}\) into \(t_{2} \otimes t_{1}\) .

13.4.2 (“Associativity and functoriality”). Let \(X_{i}\) (i = 1, 2, 3) be manifolds of class \(C^{r}\) , and let \(x_{i} \in X_{i}\) , \(t_{i} \in \mathbf{T}_{x_{i}}^{(k_{i})}(\mathbf{X}_{i})\) with \(k_{1} + k_{2} + k_{3} \leqslant r\) . One has

\[
\left(t _ {1} \otimes t _ {2}\right) \otimes t _ {3} = t _ {1} \otimes \left(t _ {2} \otimes t _ {3}\right);
\]

this point distribution is denoted \(t_{1} \otimes t_{2} \otimes t_{3}\) . One defines in the same way arbitrary finite tensor products.

Let \(\varphi_{1}\colon\mathbf{X}_{1}\to\mathbf{Y}_{1}\) and \(\varphi_{2}\colon\mathbf{X}_{2}\to\mathbf{Y}_{2}\) be morphisms of manifolds of class \(\mathbf{C}^{r}\) , and let \(t_{1}\in\mathbf{T}^{(k_{1})}(\mathbf{X}_{1})\) , \(t_{2}\in\mathbf{T}^{(k_{2})}(\mathbf{X}_{2})\) with \(k_{1}+k_{2}\leqslant r\) . One has

\[
(\varphi_ {1} \times \varphi_ {2}) _ {*} (t _ {1} \otimes t _ {2}) = \varphi_ {1 *} (t _ {1}) \otimes \varphi_ {2 *} (t _ {2}).
\]

13.4.3. With the notation of 13.4.1, let \(F_{1}\) , \(F_{2}\) and F be separated polynormed spaces, and let \((u_{1}, u_{2}) \mapsto u_{1}.u_{2}\) be a continuous bilinear map from \(F_{1} \times F_{2}\) to F. Let

\[
f _ {i} \colon \mathrm{X} _ {i} \rightarrow \mathrm{F} _ {i} \quad (i = 1, 2)
\]

be maps of class \(\mathbf{C}^r\) , and define \(f_1 \otimes f_2: \mathbf{X} \to \mathbf{F}\) by

\[
(f _ {1} \otimes f _ {2}) (y _ {1}, y _ {2}) = f _ {1} (y _ {1}). f _ {2} (y _ {2}).
\]

The map \(f_{1} \otimes f_{2}\) is of class \(\mathbf{C}^{r}\) , and we have

\[
\langle f _ {1} \otimes f _ {2}, t _ {1} \otimes t _ {2} \rangle = \langle f _ {1}, t _ {1} \rangle . \langle f _ {2}, t _ {2} \rangle .
\]

13.4.4. With the notation as in 13.4.1, let \(f\) be a map of class \(\mathbf{C}^{r}\) from X into a separated polynormed space F. If \(y_{1} \in \mathbf{X}_{1}\) , denote by \(f_{y_{1}}\) the map \(y_{2} \mapsto f(y_{1}, y_{2})\) from \(\mathbf{X}_{2}\) to F and set \(g(y_{1}) = \langle f_{y_{1}}, t_{2} \rangle\) . The function \(g: \mathbf{X}_{1} \to \mathbf{F}\) thus defined is of class \(\mathbf{C}^{r - k_{1}}\) , and we have

\[
\langle f, t _ {1} \otimes t _ {2} \rangle = \langle g, t _ {1} \rangle ,
\]

in other words

\[
\langle f, t _ {1} \otimes t _ {2} \rangle = \langle y _ {1} \mapsto \langle y _ {2} \mapsto f (y _ {1}, y _ {2}), t _ {2} \rangle , t _ {1} \rangle .
\]

We likewise have

\[
\langle f, t _ {1} \otimes t _ {2} \rangle = \langle y _ {2} \mapsto \langle y _ {1} \mapsto f (y _ {1}, y _ {2}), t _ {1} \rangle , t _ {2} \rangle .
\]

13.4.5. With the notation as in 13.4.1, let

\[
\alpha_ {k _ {1}, k _ {2}} \colon \mathrm{T} _ {x _ {1}} ^ {(k _ {1})} (\mathrm{X} _ {1}) \otimes \mathrm{T} _ {x _ {2}} ^ {(k _ {2})} (\mathrm{X} _ {2}) \rightarrow \mathrm{T} _ {x} ^ {(k)} (\mathrm{X})
\]

the linear map defined by the bilinear map \((t_{1}, t_{2}) \mapsto t_{1} \otimes t_{2}\) ; this map is injective; it is bijective if \(k_{1}\) and \(k_{2}\) are infinite.

If \(n \leqslant r\) , denote by \(\mathrm{T}_{x}^{(n)}(\mathrm{X}_{1}, \mathrm{X}_{2})\) the vector subspace of \(\mathrm{T}_{x_{1}}^{(n)}(\mathrm{X}_{1}) \otimes \mathrm{T}_{x_{2}}^{(n)}(\mathrm{X}_{2})\) generated by the subspaces \(\mathrm{T}_{x_{1}}^{(k_{1})}(\mathrm{X}_{1}) \otimes \mathrm{T}_{x_{2}}^{(k_{2})}(\mathrm{X}_{2})\) for \(k_{1} + k_{2} \leqslant n\) . There exists a linear map

\[
\alpha_ {n} \colon \mathrm{T} _ {x} ^ {(n)} (\mathrm{X} _ {1}, \mathrm{X} _ {2}) \rightarrow \mathrm{T} _ {x} ^ {(n)} (\mathrm{X})
\]

and exactly one that extends the mappings \(\alpha_{k_1,k_2}\) defined above; it is an isomorphism. In what follows, we identify \(\mathrm{T}_x^{(n)}(\mathrm{X})\) by means of \(\alpha_n^{-1}\) to the vector subspace \(\mathrm{T}_x^{(n)}(\mathrm{X}_1,\mathrm{X}_2)\) of \(\mathrm{T}_{x_1}^{(n)}(\mathrm{X}_1)\otimes \mathrm{T}_{x_2}^{(n)}(\mathrm{X}_2)\) .

13.4.6. We keep the notations of 13.4.5 and 13.3.1. By passing to the quotient, \(\alpha_{k_1,k_2}\) defines a linear map

\[
\varepsilon_ {k _ {1}, k _ {2}} \colon \operatorname{gr} _ {k _ {1}} \mathrm{T} _ {x _ {1}} ^ {(r)} (\mathrm{X} _ {1}) \otimes \operatorname{gr} _ {k _ {2}} \mathrm{T} _ {x _ {2}} ^ {(r)} (\mathrm{X} _ {2}) \rightarrow \operatorname{gr} _ {k} \mathrm{T} _ {x} ^ {(r)} (\mathrm{X}).
\]

The maps \(\varepsilon_{k_1,k_2}\) ( \(k_1 + k_2 \leqslant r\) ) are the components of a graded linear map of degree zero

\[
\varepsilon : \bigoplus_ {k _ {1} + k _ {2} \leqslant r} \left(\operatorname{gr} _ {k _ {1}} \mathrm{T} _ {x _ {1}} ^ {(r)} (\mathrm{X} _ {1}) \otimes \operatorname{gr} _ {k _ {2}} \mathrm{T} _ {x _ {2}} ^ {(r)} (\mathrm{X} _ {2})\right)\rightarrow \bigoplus_ {k \leqslant r} \operatorname{gr} _ {k} \mathrm{T} _ {x} ^ {(r)} (\mathrm{X})
\]

which is an isomorphism. The diagram

\[
\begin{array}{l} \bigoplus_ {k _ {1} + k _ {2} \leqslant r} (\operatorname{gr} _ {k _ {1}} \mathrm{T} _ {x _ {1}} ^ {(r)} (\mathrm{X} _ {1}) \otimes \operatorname{gr} _ {k _ {2}} \mathrm{T} _ {x _ {2}} ^ {(r)} (\mathrm{X} _ {2})) \xrightarrow {\varepsilon} \operatorname{gr} \mathrm{T} _ {x} ^ {(r)} (\mathrm{X}) \\ \Biggl \downarrow i _ {1 2} \\ \bigoplus_ {k _ {1} + k _ {2} \leqslant r} (\mathrm{TS} ^ {k _ {1}} (\mathrm{T} _ {x _ {1}} (\mathrm{X} _ {1})) \otimes \mathrm{TS} ^ {k _ {2}} (\mathrm{T} _ {x _ {2}} (\mathrm{X} _ {2}))) \xrightarrow {\sigma} \bigoplus_ {k \leqslant r} \mathrm{TS} ^ {k} (\mathrm{T} _ {x _ {1}} (\mathrm{X} _ {1}) \times \mathrm{T} _ {x _ {2}} (\mathrm{X} _ {2})) \end{array}
\]

is commutative (in this diagram, \(i_{12}\) denotes the homomorphism induced by \(i_{\mathbf{x}_1} \otimes i_{\mathbf{x}_2}\) and \(\sigma\) is the isomorphism defined in A, IV, § 5, no. 5). If \(r\) is infinite, this diagram is simply written:

\[
\begin{array}{l} \operatorname{gr} \mathrm{T} _ {x _ {1}} ^ {(\infty)} (\mathrm{X} _ {1}) \otimes \operatorname{gr} \mathrm{T} _ {x _ {2}} ^ {(\infty)} (\mathrm{X} _ {2}) \xrightarrow {\varepsilon} \operatorname{gr} \mathrm{T} _ {x} ^ {(\infty)} (\mathrm{X}) \\ \Biggl \downarrow_ {i _ {\mathbf {X} _ {1}} \otimes i _ {\mathbf {X} _ {2}}} \\ \operatorname{TS} (\mathrm{T} _ {x _ {1}} (\mathrm{X} _ {1})) \otimes \operatorname{TS} (\mathrm{T} _ {x _ {2}} (\mathrm{X} _ {2})) \xrightarrow {\sigma} \operatorname{TS} (\mathrm{T} _ {x} (\mathrm{X})). \end{array}
\]

